\section{Omitted Results}
\label{app:omitted}

\subsection{Convexity of the fusion losses}
\label{app:fusion-convexity}

Step 2 of Algorithm~\ref{alg:cate-fusion-selection} minimizes the losses of Definition~\ref{def:fusion-losses} over $\zeta\in\mathcal Z$. The following expansion shows that this is a convex least-squares problem. It is recorded here because the results of Section~\ref{sec:cate-risk-selection} condition on the fitted function and do not use it.

\begin{proposition}[Convexity of the fusion losses]
\label{prop:fusion-convexity}
Fix $j\in\{\mathrm{DRF},\mathrm{RF}\}$, $\lambda\in[0,1]$, $\omega\in[0,1]$, and $r\geq0$. Expanding the squares in Definition~\ref{def:fusion-losses}, $\widehat L_j(\zeta;\lambda,\omega,r)$ is a quadratic functional of $\zeta$ whose second-order part is
\begin{align}
\text{DRF: }&(1-\omega)\,\mathbb P_R^{\mathrm{tune}}\{\zeta(X)^2\}+\omega\,\mathbb P_O^{\mathrm{tune}}\{r(X)\zeta(X)^2\},
\notag\\
\text{RF: }&(1-\omega)\,\mathbb P_R^{\mathrm{tune}}\{\chi(A,X)\zeta(X)^2\}+\omega\,\mathbb P_O^{\mathrm{tune}}\{r(X)\zeta(X)^2\}.
\label{eq:fusion-curvature}
\end{align}
Both are nonnegative, so $\widehat L_j(\cdot;\lambda,\omega,r)$ is convex on any convex $\mathcal Z$. The expected loss $\mathbb E\{\widehat L_j(\zeta;\lambda,\omega,r)\}$ has the same second-order part with $\mathbb E_R$ and $\mathbb E_O$ in place of $\mathbb P_R^{\mathrm{tune}}$ and $\mathbb P_O^{\mathrm{tune}}$, and is convex as well.
\end{proposition}
Convexity is the condition under which, by Definition~\ref{def:orthogonality}, the minimizers of the expected losses are exactly the zeros of the learning equations in Definition~\ref{def:learning-equations}.

\subsection{Closed-form noise-floor coefficients}
\label{app:noise-floor}

A closed-form choice of $\lambda$ minimizes the noise floors of Proposition~\ref{prop:cate-target} instead of the CATE risk of the fitted learner in Definition~\ref{def:fitted-fusion-learner}. It is recorded here because the validation route of Section~\ref{sec:cate-risk-selection} targets the fitted risk directly and also chooses $\omega$.

\begin{proposition}[Noise-floor coefficients]
\label{prop:cate-noise-floor-lambda}
Suppose Assumptions~\ref{ass:trial} and~\ref{ass:honest-separation} hold and the displayed moments are finite. Recall $Z_0$, $Z_1$, and $\Delta=Z_1-Z_0$ from Section~\ref{sec:rct-obs-setting}, so that $Z_\lambda=Z_0+\lambda\Delta$, and write $d_m\triangleq m_O-m_R$, so that $m_\lambda=m_R+\lambda d_m$. Let $\operatorname{Proj}_{[0,1]}(u)\triangleq\min\{1,\max(0,u)\}$. If $\mathbb E_R(\Delta^2)>0$, then
\begin{align}
\operatorname*{arg\,min}_{\lambda\in[0,1]}Q_{\mathrm{DR}}(\lambda)
&=\operatorname{Proj}_{[0,1]}\left\{-\frac{\mathbb E_R(Z_0\Delta)}{\mathbb E_R(\Delta^2)}\right\},
\label{eq:cate-dr-noise-lambda}
\end{align}
and if $\mathbb E_R[d_m(X)^2/\{e(X)(1-e(X))\}]>0$, then
\begin{align}
\operatorname*{arg\,min}_{\lambda\in[0,1]}Q_R(\lambda)
&=\operatorname{Proj}_{[0,1]}\left[
\frac{\mathbb E_R[\{Y-m_R(X)\}d_m(X)/\{e(X)(1-e(X))\}]}
{\mathbb E_R[d_m(X)^2/\{e(X)(1-e(X))\}]}\right].
\label{eq:cate-r-noise-lambda}
\end{align}
If a denominator is zero, the corresponding noise floor is constant in $\lambda$. Both minimizers can be estimated by replacing the moments with tuning-sample averages.
\end{proposition}
